\scanpage{9}
\mspage{7}

(2.2.2), et on ne se bornera pas \deleted{\unsure{et}} à étudier sa restriction à $I$. \deleted{\unsure{Pour}} On verra que \deleted{toute} \inserted{parfois} le fait qu'une représentation de $I$ se prolonge à $\pi$ impose à cette représentation des restrictions draconiennes.

2.3.

2.4. Rappelons la structure de $I$. Soit $p$ l'exposant caractéristique de $k$, \deleted{Alors} et soit $P$ \deleted{\unsure{le}} un \deleted{plus grand} sous-groupe de Silow de $I$ ($P=1$ si $p=1$ i.e. $k$ de car. nulle). Alors $P$ \uline{est invariant}, i.e. on a une suite exacte
\diagram{\begin{tikzcd}[ampersand replacement=\&]
(2.4.1)\quad1\arrow[r,no head]\&P\arrow[r]\&I\arrow[r]\&I_t\arrow[r]\&1,
\end{tikzcd}}
où $I_t$ \inserted{(partie ``tame'' du groupe d'inertie $I$)} est un groupe profini premier à $p$, d'autre part on a un isomorphisme canonique
\[
(2.4.2)\qquad I_t\simeq\prod_{\ell\ne p}\mathbb Z_\ell(1)
\]
Le deuxième est obtenu ainsi. Le groupe $P$ est le groupe de Galois \deleted{de}
\[
(2.4.3)\qquad P=\operatorname{Gal}(\bar K/\bar K_t),
\]
où $\bar K_t$ est la plus \deleted{\unsure{grande}} grde sous-extension \uline{modérément ramifiée} de $\bar K/K$ \inserted{(val. : $V$)}, ou encore de $\bar K/\widetilde K$ (val. : $\widetilde V$). Celle-ci s'obtient, \unsure{en termes|en donnant}
\margin{8}
